\documentclass[11pt]{article}
\usepackage[a4paper, top=18mm, bottom=18mm, left=20mm, right=20mm]{geometry}
\usepackage[T2A]{fontenc}
\usepackage[utf8]{inputenc}
\usepackage[ukrainian]{babel}
\usepackage{graphicx}
\setlength{\parindent}{0pt}
\setlength{\parskip}{4pt}
\pagestyle{empty}
\begin{document}
%begin
{\large\textbf{open-15}}

\textbf{Варіант \No\,1}\hfill \textit{ПІБ:}\,\underline{\hspace{60mm}}
\vspace{4mm}

%beginitem
1. Що буде виведено за виконання фрагменту коду?
Записати ВИРАЗ, значенням якого є список, що складається з цілих чисел від 13 до 2003 (включно), причому спочатку за зростанням записано всі, що не діляться на 5, а потім решту за зростанням.

%enditem
%beginitem
2. Що буде виведено за виконання фрагменту коду?

%code
print('R', end=' ')
k=2
while k<8:
    k+=1
print(k)
%code


%enditem
%beginitem
3. Що буде виведено за виконання фрагменту коду?
Записати ВИРАЗ, значенням якого є список, що складається з квадратівцілих чисел від 13 до 2003 (включно), що не діляться на 5, записаних в оберненому порядку.
%enditem
%beginitem
4. Що буде виведено за виконання фрагменту коду?
True>"4"
%enditem
%beginitem
5. Що буде виведено за виконання фрагменту коду?
Записати ВИРАЗ, значенням якого є список, що складається з квадратівцілих чисел від 13 до 2003 (включно), що не діляться на 5, записаних в обернено\-му порядку.

%enditem
%beginitem
6. Що буде виведено за виконання фрагменту коду?
%code

x,y,z=2,3,0
z,x,x,z=x,8,z,y
print(x,y,z)
%code

%enditem
%beginitem
7. Що буде виведено за виконання фрагменту коду?

%code
a,b,c=2,9,0
def f(a):
    a=4
    b=3
    c=5
    return a+b+c

a,b,c=5,4,6
print(f(a),a,b,c)
%code

%enditem
%beginitem
8. Що буде виведено за виконання фрагменту коду?
try:
    k = 8
    while k>2:
        k += -3
    print(k, end=';')
except:
    print('ERR')

%enditem
%beginitem
9. Що буде виведено за виконання фрагменту коду?
k=8
   while k>2: k+=-3
   print(k)
%enditem
%beginitem
10. Що буде виведено за виконання фрагменту коду?

%code
True>"4"
%code

%enditem
%beginitem
11. Що буде виведено за виконання фрагменту коду?

%code
print('R', end=' ')
c=[10,11,12,13,14]
c[-7:-7]='11'
print(c)
%code


%enditem
%beginitem
12. Що буде виведено за виконання фрагменту коду?
-2**-2+1
%enditem
%beginitem
13. Що буде виведено за виконання фрагменту коду?

%code
def f(a,b):
    c=72
    if a>-2:
        return 2
    if b>=4:
         c=9
    else: 
        c=0
    return c

print(f(-6,-6))
%code

%enditem
%beginitem
14. Що буде виведено за виконання фрагменту коду?
try:
    for c in range(-12, -12, 2):
        if c>5:
            continue
        print(c, end=';')
        if c>5:
            break
    else:
        print(c,end=';')
    print(c)
except:
    print('Z')

%enditem
%beginitem
15. Що буде виведено за виконання фрагменту коду?

%code
6/6//6*6
%code

%enditem















%end
\newpage
\end{document}

